\part{快速开始}

\section{安装与快速开始}
\subsection{安装宏包}
支持环境为 LaTeX 2025-11-01、expl3 2025-12-29、PGF/\manualTikZ{} 3.1.11a 或更新版本。
较旧发行版应先更新，再安装本版宏包。
\pkg{gradbars} 已被 \href{https://ctan.org/pkg/gradbars}{CTAN} 收录，并以 \pkg{gradbars} 包名进入 TeX Live。
使用由 tlmgr 管理的 TeX Live（包括 MacTeX）时，可在终端安装：
\begin{lstlisting}[style=manual,language={}]
tlmgr install gradbars
\end{lstlisting}
已安装时可更新，或查看宏包信息：
\begin{lstlisting}[style=manual,language={}]
tlmgr update gradbars
tlmgr info gradbars
\end{lstlisting}
安装后在导言区写入 \key{\usepackage{gradbars}}，无需把 sty 文件复制到项目。
tlmgr 提供所配置 TeX Live 仓库中的版本，新 CTAN 版本可能尚未同步；使用本手册的新功能前请核对版本。
需要使用尚未同步的发行包时，应区分开发仓库与 CTAN 源码包。
开发仓库的 \key{sty/} 已包含运行文件；CTAN 源码包只提供可提取的 \key{.dtx/.ins}，
先在源码包根目录运行：
\begin{lstlisting}[style=manual,language={}]
texlua scripts/unpack.lua
\end{lstlisting}
该命令自动创建缺失的 \key{sty/} 并提取全部二十三个运行文件，无需 shell escape。
直接运行 \key{tex gradbars.ins} 时须确保 \key{sty/} 目录及可读的 \key{sty/README.md} 标记同时存在；
缺少任一项均使用上述 Lua 入口修复，仅创建目录不足以直接提取。
随后把 \key{sty/} 目录下的全部 \key{.sty} 文件一起放在主文档旁边，
包括可选的 \key{tblrlibgradbars.sty} 适配库。
项目中的同名文件会优先于系统安装版加载。参见 \href{https://tug.org/texlive/doc/tlmgr.html}{TeX Live 包管理说明}。
无需手动创建 \key{tikzpicture}，
也无需开启 shell escape。宏包依赖 \manualTikZ{}、xparse、expl3、collcell（含 array）、booktabs、colortbl 和 longtable，不强制加载 ctex。
中文文档可自行使用 ctexart；本手册使用 XeLaTeX 编译。

\subsection{按需加载库}
在导言区通过 \key{libraries} 选项指定需要的功能库：
\begin{codeonly}
\usepackage[libraries={stack,spark,heat,csv,table}]{gradbars}
\end{codeonly}
这段示例显式加载全部五个功能库，与默认的 \key{\usepackage{gradbars}} 等效。
只需部分功能时可缩减列表，例如：
\begin{codeonly}
\usepackage[libraries={stack,spark,heat}]{gradbars}
\end{codeonly}
\key{table} 会自动加载 \key{stack}、\key{spark}、\key{heat} 和 \key{csv} 依赖；
\key{libraries=none} 只加载核心与公开接口。库加载须在导言区完成，详见后面的库加载章节。

\subsection{最小完整文档}
安装宏包后，将下面代码保存为一个 \key{.tex} 文件即可编译；手动安装时将全部运行 sty 文件放在同一目录。
\begin{codeonly}
\documentclass{article}
\usepackage{gradbars}
\begin{document}
\gradbarssetup{max=100, width=24mm, precision=1}
\gradbar{72.4}

\gradbar{92.7}
\end{document}
\end{codeonly}


\subsection{第一张可视化表格}
下面是可直接编译的完整文档。得分越高越好，耗时越低越好；缺失数据用 NA，不能用零代替。
\begin{codeonly}
\documentclass{article}
\usepackage{gradbars}
\begin{document}
\begin{center}
\begin{gradtable}[width=\linewidth,header align=c]
  \gradcolumn{method}{Method}
  \gradcolumn[type=bar,mark=both,
    options={max=100,precision=1}]{score}{Score}
  \gradcolumn[type=number,mark=best,
    options={better=lower,precision=1}]{time}{Time / ms}
  \gradrow{Baseline,72.4,18.2}
  \gradrow{Proposed,92.7,11.5}
  \gradrow{Pending,NA,NA}
\end{gradtable}
\end{center}
\end{document}
\end{codeonly}
\begin{example}[text and listing,sidebyside=false]{第一张表格的实际效果与正文代码}
\begin{center}
\begin{gradtable}[width=\linewidth,header align=c]
  \gradcolumn{method}{Method}
  \gradcolumn[type=bar,mark=both,
    options={max=100,precision=1}]{score}{Score}
  \gradcolumn[type=number,mark=best,
    options={better=lower,precision=1}]{time}{Time / ms}
  \gradrow{Baseline,72.4,18.2}
  \gradrow{Proposed,92.7,11.5}
  \gradrow{Pending,NA,NA}
\end{gradtable}
\end{center}
\end{example}

\subsection{首次编译排错}
\begin{description}
\item[找不到 gradbars.sty] 检查安装位置；手动安装须复制全部运行模块，并保持版本一致。
\item[中文或字体报错] 中文文档使用 XeLaTeX 与 ctexart；上述英文最小文档也可使用 pdfLaTeX 或 LuaLaTeX。
\item[命令未定义] 检查是否加载命令所属功能库，所有库加载均放在导言区。
\item[Invalid value] 先检查是否把单位、分组符号或表达式写进数值参数；地区输入须显式声明规则。
\end{description}
编译成功后，先阅读共同基础，再按图形详解选型；已有普通表格可直接查阅表格实战，
使用外部数据时转到数据与报告。参数速查与常见问题位于最后一部分。
